% ===========================================================================
%  chapter3.tex  --  Глава 3. Экономическая часть (общая со стартапом «ГМ-ПВР»)
%  Источник: общая экономическая глава стартапа (Д. С. Лыков).
%  Перенесено в ВКР А. К. Лазарева практически без изменений; заменены только
%  (1) имя разработчика-исследователя и (2) перечень нейросетевых моделей.
% ===========================================================================

% --- 4-й уровень нумерации (3.1.2.1) для экономической главы ---------------
\setcounter{secnumdepth}{4}
\titleformat{\paragraph}
  {\centering\bfseries\fontsize{13pt}{15.6pt}\selectfont}
  {\theparagraph.}{0.5em}{}
\titlespacing*{\paragraph}{0pt}{12pt}{6pt}

% --- Вспомогательные типы колонок ------------------------------------------
\newcolumntype{L}[1]{>{\raggedright\arraybackslash}p{#1}}
\newcolumntype{C}[1]{>{\centering\arraybackslash}p{#1}}

% --- Команда для ручной подписи таблиц (нумерация как в задании) ------------
\newcommand{\tcap}[1]{\par\vspace{4pt}\noindent{\itshape\bfseries #1}\par\nopagebreak\vspace{2pt}}

\section{Экономическая часть}
\label{sec:econom}

\subsection{Экономическое обоснование проекта}

В~выполнении работ принимают участие два специалиста: руководитель проекта~---
Петрусевич~Д.\,А., к.\,ф.-м.\,н., доцент, и~разработчик-исследователь~---
студент 4-го курса Лазарев~А.\,К., группа КМБО-03-22.
Руководитель проекта отвечает за~корректную формулировку исследовательской
задачи, контроль методологии, проверку обоснованности используемого
математического аппарата при~описании вейвлет-разложений и~нейросетевых
архитектур, вносит необходимые коррективы и~оценивает выполненную работу
в~целом.
Разработчик-исследователь выполняет основные этапы проекта: проектирует
архитектуру решения, проводит анализ данных на~синтетических рядах и~наборах
M3/M4, реализует декомпозицию временного ряда на~компоненты, разрабатывает
и~исследует нейросетевые модели, применяемые к~компонентам временного ряда,
полученным с~помощью вейвлет-декомпозиции, реализует REST~API, выполняет
тестирование системы и~подготавливает проектную документацию.

\subsubsection{Организация и планирование работ по~теме}

На~работу отводится 90~рабочих дней. Этапы разработки, включающие
исследовательскую часть и~реализацию с~использованием выбранного программного
обеспечения (Python, PyTorch Forecasting, Sktime, FastAPI, PostgreSQL),
представлены в~таблице~3.1.1.1.

\tcap{Таблица 3.1.1.1. Этапы разработки стартапа «ГМ-ПВР».}
{\small
\begin{longtable}{C{0.7cm} L{9.0cm} L{2.4cm} C{2.0cm}}
\toprule
№ & Название этапа & Исполнитель & Трудоёмкость, чел.-дни \\
\midrule
\endhead
\multirow{2}{*}{1} & \multirow{2}{9.0cm}{Анализ предметной области и~обзор существующих подходов к~краткосрочному прогнозированию временных рядов} & Разработчик & 12 \\*
 & & Руководитель & 3 \\
\midrule
\multirow{2}{*}{2} & \multirow{2}{9.0cm}{Изучение методов декомпозиции временных рядов: STL, DWT, MODWT, EWT и~VMD} & Разработчик & 10 \\*
 & & Руководитель & 3 \\
\midrule
\multirow{2}{*}{3} & \multirow{2}{9.0cm}{Разработка архитектурной концепции гибридной модели HybridPlus на~основе MODWT/STL-декомпозиции} & Разработчик & 12 \\*
 & & Руководитель & 3 \\
\midrule
\multirow{2}{*}{4} & \multirow{2}{9.0cm}{Реализация базовых статистических моделей ARIMA, ARIMA с~автоподбором, ETS и~Prophet} & Разработчик & 10 \\*
 & & Руководитель & 2 \\
\midrule
\multirow{2}{*}{5} & \multirow{2}{9.0cm}{Реализация нейросетевых моделей прогнозирования компонент ряда: N-BEATS, DLinear, TimesNet, Autoformer, FEDformer, PatchTST, Helformer} & Разработчик & 12 \\*
 & & Руководитель & 3 \\
\midrule
\multirow{2}{*}{6} & \multirow{2}{9.0cm}{Разработка гибридных схем прогнозирования: MODWT~+ нейросетевые модели~+ ARIMA/ETS и~STL-гибридов} & Разработчик & 14 \\*
 & & Руководитель & 4 \\
\midrule
\multirow{2}{*}{7} & \multirow{2}{9.0cm}{Проведение экспериментов на~синтетических рядах, наборах M3/M4 и~рядах с~многопериодной сезонностью и~сложным трендом} & Разработчик & 12 \\*
 & & Руководитель & 3 \\
\midrule
\multirow{2}{*}{8} & \multirow{2}{9.0cm}{Реализация REST~API сервиса и~подготовка структуры хранения результатов экспериментов} & Разработчик & 8 \\*
 & & Руководитель & 2 \\
\midrule
\multirow{2}{*}{9} & \multirow{2}{9.0cm}{Анализ результатов, экономическое обоснование, оформление ВКР и~подготовка презентации} & Разработчик & 10 \\*
 & & Руководитель & 2 \\
\midrule
\multicolumn{3}{r}{\textbf{Итого, разработчик}} & \textbf{90} \\
\multicolumn{3}{r}{\textbf{Итого, руководитель}} & \textbf{25} \\
\bottomrule
\end{longtable}
}

\subsubsection{Расчёт себестоимости проекта}

Расчёт себестоимости проекта с~учётом затрат на~специальное оборудование,
заработную плату всех участников, дополнительную заработную плату, единый
страховой платёж с~учётом взноса на~травматизм, накладные и~прочие расходы
в~ВКР был проведён методом калькулирования~\cite{gavrilenko2019methodical,nazarova2021ekonomika}.

\paragraph{Материалы, покупные изделия и полуфабрикаты}

В~ходе выполнения ВКР возникла необходимость приобретения материалов, изделий
и~иных материальных ресурсов, необходимых для~проведения исследования
и~оформления результатов работы. В~данную статью затрат также включены расходы,
связанные с~подготовкой комплектов документации. Кроме того, при~расчёте
материальных затрат учитываются транспортно-заготовительные расходы. Состав
и~стоимость материальных затрат представлены в~таблице~3.1.2.1.

\tcap{Таблица 3.1.2.1. Стоимость затрат на~материальные ценности.}
{\small
\begin{longtable}{C{0.8cm} L{5.6cm} C{2.2cm} C{1.8cm} C{2.4cm} C{2.4cm}}
\toprule
№ пп & Наименование материалов & Единицы измерения & Количество & Цена за единицу, руб. & Стоимость, руб. \\
\midrule
\endhead
1 & Бумага А4 & пачка & 1 & 399 & 399 \\
2 & Картридж для~принтера & шт & 1 & 3711 & 3711 \\
3 & Ручка & шт & 10 & 50 & 500 \\
4 & Карандаш & шт & 5 & 20 & 100 \\
\midrule
\multicolumn{5}{r}{Итого материалов} & 4710 \\
\multicolumn{5}{r}{Транспортно-заготовительные расходы} & 500 \\
\multicolumn{5}{r}{\textbf{Итого}} & \textbf{5210} \\
\bottomrule
\end{longtable}
}

\paragraph{Специальное оборудование}

Ввиду использования тяжёлых нейросетевых моделей (трансформерных архитектур
Autoformer, FEDformer, PatchTST, Helformer, а~также N-BEATS, DLinear и~TimesNet),
а~также необходимости проведения большого числа экспериментов на~наборах
временных рядов, в~проекте учитываются расходы на~аренду облачных
вычислительных ресурсов. Для~хранения промежуточных результатов и~логов
экспериментов дополнительно учитывается облачное хранилище. Расчёт затрат
на~специальное оборудование приведён в~таблице~3.1.2.2.

\tcap{Таблица 3.1.2.2. Стоимость затрат на~специальное оборудование.}
{\small
\begin{longtable}{C{0.8cm} L{5.6cm} C{2.2cm} C{1.8cm} C{2.4cm} C{2.4cm}}
\toprule
№ пп & Наименование & Единицы измерения & Количество & Цена за единицу, руб. & Стоимость, руб. \\
\midrule
\endhead
1 & Аренда облачного GPU-сервера & мес. & 3 & 8000 & 24000 \\
2 & Облачное хранилище для~результатов экспериментов & мес. & 3 & 1000 & 3000 \\
\midrule
\multicolumn{5}{r}{\textbf{Итого}} & \textbf{27000} \\
\bottomrule
\end{longtable}
}

\paragraph{Основная заработная плата}

В~таблице~3.1.2.3 представлен расчёт основной заработной платы участников
проекта. Расчёт выполнен на~основе трудоёмкости этапов разработки и~дневной
ставки исполнителей, определяемой исходя из~месячного оклада (месячный оклад
руководителя~--- 60\,000~руб./мес., разработчика~--- 50\,000~руб./мес.).

\tcap{Таблица 3.1.2.3. Расчёт основной заработной платы.}
{\footnotesize
\begin{longtable}{C{0.5cm} L{4.7cm} L{1.9cm} C{1.4cm} C{1.3cm} C{1.4cm} C{1.6cm}}
\toprule
№ & Наименование этапа & Исполнитель & Мес. оклад, руб. & Труд., чел.-дни & Оплата за день, руб. & Оплата за этап, руб. \\
\midrule
\endhead
\multirow{2}{*}{1} & \multirow{2}{4.7cm}{Анализ предметной области и~обзор существующих подходов} & Руководитель & 60\,000 & 3 & 2\,408 & 7\,224 \\*
 & & Разработчик & 50\,000 & 12 & 2\,007 & 24\,084 \\
\midrule
\multirow{2}{*}{2} & \multirow{2}{4.7cm}{Изучение методов декомпозиции: STL, DWT, MODWT, EWT, VMD} & Руководитель & 60\,000 & 3 & 2\,408 & 7\,224 \\*
 & & Разработчик & 50\,000 & 10 & 2\,007 & 20\,070 \\
\midrule
\multirow{2}{*}{3} & \multirow{2}{4.7cm}{Архитектурная концепция HybridPlus (MODWT/STL)} & Руководитель & 60\,000 & 3 & 2\,408 & 7\,224 \\*
 & & Разработчик & 50\,000 & 12 & 2\,007 & 24\,084 \\
\midrule
\multirow{2}{*}{4} & \multirow{2}{4.7cm}{Базовые статистические модели ARIMA, ETS, Prophet} & Руководитель & 60\,000 & 2 & 2\,408 & 4\,816 \\*
 & & Разработчик & 50\,000 & 10 & 2\,007 & 20\,070 \\
\midrule
\multirow{2}{*}{5} & \multirow{2}{4.7cm}{Нейросетевые модели: N-BEATS, DLinear, TimesNet, Autoformer, FEDformer, PatchTST, Helformer} & Руководитель & 60\,000 & 3 & 2\,408 & 7\,224 \\*
 & & Разработчик & 50\,000 & 12 & 2\,007 & 24\,084 \\
\midrule
\multirow{2}{*}{6} & \multirow{2}{4.7cm}{Гибридные схемы: MODWT~+ нейросеть~+ ARIMA/ETS и~STL} & Руководитель & 60\,000 & 4 & 2\,408 & 9\,632 \\*
 & & Разработчик & 50\,000 & 14 & 2\,007 & 28\,098 \\
\midrule
\multirow{2}{*}{7} & \multirow{2}{4.7cm}{Эксперименты на~M3/M4 и~рядах со~сложной сезонностью} & Руководитель & 60\,000 & 3 & 2\,408 & 7\,224 \\*
 & & Разработчик & 50\,000 & 12 & 2\,007 & 24\,084 \\
\midrule
\multirow{2}{*}{8} & \multirow{2}{4.7cm}{Реализация REST~API и~хранения результатов} & Руководитель & 60\,000 & 2 & 2\,408 & 4\,816 \\*
 & & Разработчик & 50\,000 & 8 & 2\,007 & 16\,056 \\
\midrule
\multirow{2}{*}{9} & \multirow{2}{4.7cm}{Анализ результатов, оформление ВКР и~презентации} & Руководитель & 60\,000 & 2 & 2\,408 & 4\,816 \\*
 & & Разработчик & 50\,000 & 10 & 2\,007 & 20\,070 \\
\midrule
\multicolumn{6}{r}{\textbf{Итого}} & \textbf{240\,830} \\
\bottomrule
\end{longtable}
}

\paragraph{Дополнительная заработная плата}

К~данной статье калькуляции относятся выплаты, предусмотренные ТК~РФ
за~время, не~отработанное по~уважительным причинам, а~также оплата отпусков.
В~среднем размер дополнительной заработной платы принимается на~уровне
20\,--\,25~\% от~основной заработной платы. В~рамках данного расчёта
используется норматив 20~\% от~основной заработной платы всех участников
проекта:
\[ \text{ДЗП} = 240\,830 \times 0{,}2 = 48\,166~\text{руб.} \]

\paragraph{Страховые отчисления}

Согласно ОКВЭД, деятельность стартапа относится к~обработке данных,
предоставлению услуг по~размещению информации и~связанной с~этим деятельности,
что соответствует коду ОКВЭД 63.11. При~этом организация не~имеет
IT-аккредитации, поэтому пониженные тарифы страховых взносов не~применяются.
Страховые взносы уплачиваются единым платежом; к~базовой ставке 30~\%
дополнительно учитывается ставка страхования от~несчастных случаев
и~профессиональных заболеваний 0,2~\%. Поэтому расчёт выполняется
по~совокупной ставке 30,2~\% от~фонда оплаты труда, включающего основную
и~дополнительную заработную плату:
\[ \text{ФОТ} = \text{ОЗП} + \text{ДЗП} = 240\,830 + 48\,166 = 288\,996~\text{руб.} \]
\[ \text{СВ} = \text{ФОТ} \times 30{,}2~\% = 288\,996 \times 0{,}302 = 87\,276{,}8~\text{руб.} \]

\paragraph{Командировочные расходы}

Расходы по~данному разделу отсутствуют.

\paragraph{Контрагентские услуги}

В~процессе разработки ВКР услуги подрядных организаций не~использовались.

\paragraph{Накладные расходы}

Накладные расходы (НР) представляют собой затраты, связанные с~организацией,
управлением и~обслуживанием процесса разработки. Величина накладных расходов
определяется в~процентах от~основной заработной платы. В~зависимости
от~структуры предприятия данный показатель обычно находится в~диапазоне
200\,--\,300~\%. Для~рассматриваемого проекта принимается минимальное значение
норматива~--- 200~\% от~ОЗП. Расчёт выполняется по~формуле:
\[ \text{НР} = \text{ОЗП} \times 200~\% = 240\,830 \times 2 = 481\,660~\text{руб.} \]

\paragraph{Прочие расходы}

Расходы по~данному разделу отсутствуют.

\paragraph{Полная себестоимость проекта}

Полная себестоимость стартапа представляет собой сумму всех описанных выше
статей калькуляции, она представлена в~таблице~3.1.2.4.

\tcap{Таблица 3.1.2.4. Полная себестоимость проекта.}
{\small
\begin{longtable}{C{0.8cm} L{11.0cm} C{3.0cm}}
\toprule
№ пп & Номенклатура статей расходов & Затраты, руб. \\
\midrule
\endhead
1 & Материалы, покупные изделия и~полуфабрикаты (за~вычетом отходов) & 5\,210 \\
2 & Специальное оборудование для~научных (экспериментальных) работ & 27\,000 \\
3 & Основная заработная плата & 240\,830 \\
4 & Дополнительная заработная плата & 48\,166 \\
5 & Страховые взносы в~социальные фонды & 87\,276,8 \\
6 & Расходы на~научные и~производственные командировки & --- \\
7 & Услуги подрядных организаций & --- \\
8 & Накладные расходы & 481\,660 \\
9 & Прочие расходы & --- \\
\midrule
\multicolumn{2}{r}{\textbf{Итого}} & \textbf{890\,142,8} \\
\bottomrule
\end{longtable}
}

Полная себестоимость разработки сервиса «ГМ-ПВР» составляет 890\,142,8~руб.
Наибольшую долю в~структуре затрат занимают накладные расходы (54,1~\%)
и~основная заработная плата (27,1~\%), что~характерно для~научно-иссле\-довательских
проектов в~области машинного обучения и~анализа данных, где~главным
ресурсом является интеллектуальный труд исполнителей.

\subsubsection{Бизнес-план и коммерческое обоснование проекта}

\paragraph{Бизнес-план}

Технологический стартап «ГМ-ПВР» представляет собой B2B-сервис
для~краткосрочного прогнозирования временных рядов через~REST~API. Целевые
клиенты~--- компании малого и~среднего бизнеса в~ритейле, логистике,
энергетике и~телекоммуникациях, которым необходимо прогнозировать спрос,
продажи, складские остатки, нагрузку, трафик или~потребление ресурсов.
Проект решает проблему нестабильного качества прогнозов при~использовании
одной универсальной модели. Реальные временные ряды часто содержат тренд,
сезонность, шум, выбросы и~структурные сдвиги, поэтому ошибки прогноза могут
приводить к~избыточным запасам, дефициту товаров, неэффективному планированию
и~дополнительным затратам.
Самостоятельная разработка такого решения требует найма специалиста
по~машинному обучению, настройки моделей, проведения экспериментов
и~поддержки инфраструктуры. Стоимость содержания одного специалиста с~учётом
страховых отчислений может составлять около 325\,000~руб./мес., тогда как
готовый сервис по~подписке позволяет снизить эти затраты.
Сервис «ГМ-ПВР» принимает временной ряд и~параметры прогноза, после~чего
возвращает результат, рассчитанный гибридной моделью HybridPlus. В~основе
подхода используется MODWT-декомпозиция: низкочастотная компонента
прогнозируется нейросетевыми моделями (Helformer, Autoformer, FEDformer,
PatchTST, TimesNet, N-BEATS, DLinear), а~высокочастотные и~остаточные
компоненты~--- статистическими моделями ARIMA или~ETS.
Модель монетизации представляет собой подписку на~сервис:

\begin{itemize}[nosep]
  \item Базовый тариф: 85\,000~руб./мес. (до~50\,000 обращений к~API прогнозирования);
  \item Дополнительная плата за~каждый API-запрос сверх лимита.
\end{itemize}

В~таблице~3.1.3.1 приведено комплексное описание продукта и~его целевых
потребителей.

\tcap{Таблица 3.1.3.1. Описание продукта и~целевых потребителей.}
{\footnotesize
\begin{longtable}{L{7.6cm} L{7.6cm}}
\toprule
\textbf{Потребители} & \textbf{Продукт} \\
\midrule
\endhead
\textbf{Проблема:} компании малого и~среднего бизнеса вынуждены прогнозировать спрос, продажи, нагрузку и~потребление ресурсов в~условиях ограниченной истории, шума, сезонности и~структурных сдвигов. Универсальная модель не~обеспечивает стабильного качества прогноза, что~приводит к~избыточным запасам, дефициту товаров и~неэффективному планированию. & \textbf{Решение:} B2B-сервис «ГМ-ПВР», предоставляющий через~REST~API доступ к~гибридной модели HybridPlus на~базе MODWT-декомпозиции, в~которой низкочастотные компоненты прогнозируются нейросетевыми моделями (N-BEATS, DLinear, TimesNet, Autoformer, FEDformer, PatchTST, Helformer), а~высокочастотные и~остаточные~--- моделями ARIMA/ETS. \\
\midrule
\textbf{Кто потребители (ЦА)?} 1)~Ритейл (прогноз спроса и~продаж); 2)~логистические компании (планирование загрузки и~маршрутов); 3)~энергетика (прогноз потребления и~нагрузки); 4)~телекоммуникации (прогноз трафика); 5)~e-commerce и~маркетплейсы (планирование запасов). & \textbf{Как работает продукт?} Клиент отправляет HTTP-запрос в~REST~API с~временным рядом и~параметрами прогноза. Сервис автоматически выполняет MODWT-разложение ряда, прогнозирует компоненты соответствующими моделями (нейросетевыми или~статистическими) и~возвращает реконструированный прогноз вместе с~показателями его~устойчивости. \\
\midrule
\textbf{Как они решают проблему сейчас?} 1)~Используют простые встроенные модели в~BI-системах; 2)~содержат собственного Data Scientist (250\,000+~руб./мес. с~учётом отчислений); 3)~применяют одну универсальную модель (ARIMA или~одиночная нейросеть) без~декомпозиции; 4)~не~прогнозируют вовсе, опираясь на~экспертные оценки. & \textbf{В~чём преимущество?} 1)~Снижение ошибок прогноза за~счёт раздельного моделирования низко- и~высокочастотных компонент; 2)~подписочная модель~--- снижение совокупных издержек по~сравнению с~собственным специалистом; 3)~автоматический подбор моделей по~структуре ряда; 4)~стабилизация первого шага прогноза для~краткосрочных задач; 5)~отсутствие необходимости в~ML-инфраструктуре на~стороне клиента. \\
\bottomrule
\end{longtable}
}

Ценностное предложение сервиса «ГМ-ПВР» заключается в~предоставлении компаниям
малого и~среднего бизнеса готового инструмента краткосрочного прогнозирования
с~устойчивым качеством на~разнородных временных рядах, без~необходимости
содержать собственного специалиста по~машинному обучению и~поддерживать
ML-инфраструктуру.

\paragraph{Целевая аудитория}

Целевой аудиторией сервиса являются компании малого и~среднего бизнеса,
в~операционной деятельности которых требуется регулярное краткосрочное
прогнозирование временных рядов: продаж, спроса, трафика, нагрузки,
потребления ресурсов. В~таблице~3.1.3.2 представлен портрет потребителя
B2B-сегмента.

\tcap{Таблица 3.1.3.2. Портрет потребителя (B2B-сегмент).}
{\small
\begin{longtable}{C{0.6cm} L{4.2cm} L{9.4cm}}
\toprule
№ & Параметр & Характеристика \\
\midrule
\endhead
1 & Организационно-правовая форма & ООО, ОАО, АО, ИП \\
2 & Отрасль (ОКВЭД) & 47 (розничная торговля), 49 (транспорт), 35 (электроэнергетика), 61 (телеком), 62 (разработка ПО), 63 (обработка данных) \\
3 & Размер штата & От~20 сотрудников \\
4 & Годовая выручка & От~50 млн руб. \\
5 & Географическое расположение & Российская Федерация (потенциально~--- страны СНГ) \\
6 & Лицо, принимающее решения & Руководитель аналитического отдела, директор по~операциям, ИТ-директор \\
7 & Требования к~инфраструктуре & Возможность интеграции с~REST~API; наличие систем учёта продаж/трафика/нагрузки \\
8 & Объём данных & Временные ряды длиной от~100 наблюдений; горизонт прогноза 1\,--\,30 шагов \\
\bottomrule
\end{longtable}
}

По~данным Росстата и~аналитики РБК, в~России насчитывается порядка
40\,--\,50~тыс. компаний малого и~среднего бизнеса в~ритейле, логистике
и~e-commerce, для~которых задача краткосрочного прогнозирования временных
рядов является регулярной операционной потребностью. Тренд на~цифровизацию
и~переход от~облачных зарубежных ML-сервисов к~отечественным решениям
после~2022~года формирует устойчивый спрос на~сервисы подобного класса.

\paragraph{SWOT-анализ}

SWOT-анализ~--- метод стратегического планирования, позволяющий выявить
сильные и~слабые стороны проекта, а~также возможности и~угрозы внешней среды.
В~таблице~3.1.3.3 представлен SWOT-анализ проекта «ГМ-ПВР».

\tcap{Таблица 3.1.3.3. SWOT-анализ проекта.}
{\footnotesize
\begin{longtable}{L{7.6cm} L{7.6cm}}
\toprule
\textbf{Сильные стороны (Strengths)} & \textbf{Слабые стороны (Weaknesses)} \\
\midrule
\endhead
1)~Гибридная архитектура HybridPlus с~раздельным моделированием низко-
и~высокочастотных компонент. 2)~Экспериментально подтверждённое преимущество
MODWT-гибридов на~рядах со~сложной сезонностью и~многошкальной структурой.
3)~Использование открытого стека: Python, PyTorch Forecasting, Sktime,
FastAPI, PostgreSQL. 4)~Стабилизация первого шага прогноза, важная
для~краткосрочных задач. 5)~Гибкая подписочная модель монетизации
с~низким порогом входа.
&
1)~На~рядах малой длины (yearly, monthly из~M4) преимущество MODWT-гибридов
нивелируется. 2)~На~часовых данных декомпозиция может приводить к~фазовым
искажениям. 3)~Зависимость от~внешних поставщиков GPU-вычислений. 4)~Узкая
специализация~--- только унивариативные ряды. 5)~Отсутствие графического
интерфейса (только API). \\
\midrule
\textbf{Возможности (Opportunities)} & \textbf{Угрозы (Threats)} \\
\midrule
А)~Курс на~импортозамещение и~переход от~зарубежных ML-сервисов
к~отечественным. Б)~Расширение функционала на~мультивариативные ряды
и~работу с~экзогенными признаками. В)~Подключение дополнительных декомпозиций
(EWT, VMD) для~специфичных профилей. Г)~Партнёрства с~поставщиками BI-систем
(Yandex DataLens, PIX BI). Д)~Выход на~рынки СНГ и~стран БРИКС.
&
1)~Конкуренция с~зарубежными облачными ML-платформами (Amazon Forecast,
Azure ML). 2)~Появление open-source альтернатив с~сопоставимым качеством.
3)~Риск разработки внутренней системы силами клиента (in-house Data
Scientist). 4)~Изменение тарифов на~GPU-аренду со~стороны облачных
провайдеров. 5)~Регуляторные ограничения на~трансграничную передачу данных. \\
\bottomrule
\end{longtable}
}

По~результатам анализа ключевым конкурентным преимуществом является сочетание
экспериментально подтверждённой эффективности гибридной MODWT-декомпозиции
на~сложных профилях и~доступной подписочной модели для~российских компаний
малого и~среднего бизнеса. Основные риски связаны с~конкуренцией крупных
облачных платформ и~появлением open-source решений, однако они могут быть
снижены за~счёт акцентирования внимания на~устойчивости прогноза
на~проблемных профилях рядов и~гибкой ценовой политики.

\paragraph{Анализ конкурентов}

На~рынке услуг прогнозирования временных рядов через~API можно выделить
несколько типов конкурентов: глобальные облачные ML-платформы,
специализированные сервисы прогнозирования и~open-source библиотеки,
на~базе которых клиент может построить собственное решение. Систематизированный
анализ конкурентов приведён в~таблице~3.1.3.4.

\tcap{Таблица 3.1.3.4. Анализ конкурентных решений.}
{\footnotesize
\begin{longtable}{L{2.5cm} L{3.9cm} L{3.6cm} L{3.3cm} C{1.3cm}}
\toprule
Решение & Подход & Преимущества & Недостатки & Страна \\
\midrule
\endhead
«ГМ-ПВР» (разрабатываемый сервис) & Гибридная модель HybridPlus с~MODWT-декомпозицией: нейросетевые модели (N-BEATS, DLinear, TimesNet, Autoformer, FEDformer, PatchTST, Helformer) для~низкочастотных компонент, ARIMA/ETS для~высокочастотных & Адаптация к~сложной сезонности и~шумному тренду; стабилизация первого шага; подписка от~85\,000~руб./мес.; российская юрисдикция & Только унивариативные ряды; на~рядах малой длины эффект MODWT нивелируется & РФ \\
\midrule
Amazon Forecast & AutoML с~выбором из~CNN-QR, DeepAR+, Prophet, ARIMA, ETS, NPTS & Высокая зрелость; интеграция с~экосистемой AWS; поддержка экзогенных признаков & Высокая стоимость (от~\$0,088 за~1000 прогнозов); санкционные ограничения для~РФ; зависимость от~AWS & США \\
\midrule
Azure ML Time Series & AutoML на~основе ансамблей (ExponentialSmoothing, Arima, Prophet, ForecastTCN) & Глубокая интеграция с~Azure; визуальные ML-конструкторы & Санкционные ограничения; высокая стоимость; сложность настройки под~нишевые задачи & США \\
\midrule
Yandex DataSphere / Yandex Cloud ML & Notebooks~+ JupyterLab с~предустановленными ML-библиотеками & Российская юрисдикция; гибкость; интеграция с~Yandex Cloud & Требуется собственный Data Scientist; не~«коробочное» решение; нет специализированных гибридных моделей & РФ \\
\midrule
Open-source библиотеки (Darts, NeuralForecast, GluonTS) & Самостоятельная разработка прогнозов на~базе библиотек & Бесплатно; полный контроль над~моделями; широкий выбор архитектур & Требуется собственная команда (Data Scientist, DevOps); затраты на~инфраструктуру; долгий time-to-market & Глоб. \\
\bottomrule
\end{longtable}
}

Анализ показал, что~основные глобальные конкуренты (Amazon Forecast, Azure ML)
являются зарубежными облачными решениями, что~после~2022~года создаёт
санкционные риски и~юридические ограничения для~большинства российских
заказчиков. Yandex DataSphere предоставляет среду для~самостоятельной
разработки, но~не~является готовым продуктом для~прогнозирования. Open-source
путь требует от~клиента содержания собственной команды специалистов. Сервис
«ГМ-ПВР» занимает нишу отечественного «коробочного» API-решения
с~экспериментально подтверждённой эффективностью на~сложных профилях рядов.

\paragraph{Бизнес-модель по~Остервальдеру}

Бизнес-модель Остервальдера (Business Model Canvas)~--- инструмент
стратегического управления для~описания и~визуализации бизнес-моделей. Она
представляет собой схему из~9~блоков, описывающих ключевые партнёры, виды
деятельности, ресурсы, ценностные предложения, взаимоотношения с~клиентами,
каналы сбыта, потребительские сегменты, структуру издержек и~потоки доходов.
В~таблицах~3.1.3.5\,--\,3.1.3.7 представлена бизнес-модель проекта «ГМ-ПВР».

\tcap{Таблица 3.1.3.5. Бизнес-модель по~Остервальдеру (Canvas)~--- часть~1.}
{\footnotesize
\begin{longtable}{L{5.0cm} L{5.0cm} L{5.0cm}}
\toprule
\textbf{Ключевые партнёры} & \textbf{Ключевые виды деятельности} & \textbf{Ценностные предложения} \\
\midrule
\endhead
Поставщики GPU-вычислений (Yandex Cloud, VK Cloud, Selectel); провайдеры PaaS-инфраструктуры; разработчики open-source библиотек (PyTorch Forecasting, Sktime); технологические партнёры из~BI-сегмента; научно-исследовательские центры по~машинному обучению. & Разработка и~сопровождение гибридной модели HybridPlus; поддержка REST~API; проведение экспериментов и~валидация на~новых профилях рядов; обновление моделей и~переобучение нейросетевых компонент; техническая поддержка клиентов; маркетинг и~продажи. & Для~ритейла и~e-commerce: снижение ошибок прогноза спроса и~потерь от~избытка/дефицита запасов на~25~\% и~более. Для~логистики и~телекома: устойчивое прогнозирование нагрузки при~сложной многопериодной сезонности. Для~всех сегментов: отсутствие необходимости в~собственном Data Scientist и~ML-инфраструктуре. \\
\bottomrule
\end{longtable}
}

\tcap{Таблица 3.1.3.6. Бизнес-модель по~Остервальдеру (Canvas)~--- часть~2.}
{\footnotesize
\begin{longtable}{L{5.0cm} L{5.0cm} L{5.0cm}}
\toprule
\textbf{Ключевые ресурсы} & \textbf{Каналы сбыта} & \textbf{Потребительские сегменты} \\
\midrule
\endhead
\textit{Кадровые:} команда из~ML-инженеров, backend-разработчиков и~DevOps. \textit{Технологические:} гибридная модель HybridPlus и~кодовая база сервиса. \textit{Финансовые:} собственные средства и~грантовое финансирование. \textit{Инфраструктурные:} облачные GPU-серверы для~обучения и~инференса. \textit{Интеллектуальные:} накопленные результаты экспериментов на~M3/M4 и~синтетических рядах. & Прямые продажи через~сайт компании; контент-маркетинг (статьи на~Habr, vc.ru, отраслевых порталах); участие в~ML-конференциях (Data Fusion, AI~Journey, OpenTalks.AI); технические демонстрации и~пилотные проекты; партнёрский канал через~BI- и~системных интеграторов. & B2B: средний и~малый бизнес в~ритейле, e-commerce, логистике, энергетике, телекоммуникациях; аналитические подразделения крупных корпораций (как~замена части in-house задач); консалтинговые компании, оказывающие услуги анализа данных. \\
\bottomrule
\end{longtable}
}

\tcap{Таблица 3.1.3.7. Бизнес-модель по~Остервальдеру (Canvas)~--- часть~3.}
{\footnotesize
\begin{longtable}{L{7.6cm} L{7.6cm}}
\toprule
\textbf{Структура издержек} & \textbf{Потоки доходов} \\
\midrule
\endhead
\textit{Постоянные:} ФОТ команды разработки и~поддержки; аренда офиса; постоянная часть платы за~облачную инфраструктуру; сопровождение домена и~сайта. \textit{Переменные:} аренда GPU-серверов по~объёму инференса; маркетинг и~реклама; командировки на~конференции; комиссии платёжных систем; налоги (УСН). & Базовая подписка 85\,000~руб./мес. (до~50\,000 API-запросов); плата за~превышение лимитов (по~тарификации за~пакеты запросов); расширенный тариф с~поддержкой кастомных профилей (от~200\,000~руб./мес.); разовые контракты на~пилотные проекты и~валидацию моделей под~конкретного клиента (от~300\,000~руб.). \\
\bottomrule
\end{longtable}
}

Основным потоком доходов является подписочная плата за~использование REST~API.
Подписочная модель обеспечивает прогнозируемый денежный поток и~масштабируется
без~существенного увеличения постоянных издержек, поскольку дополнительный
клиент потребляет преимущественно облачные GPU-ресурсы, расходы на~которые
являются переменными.

\paragraph{Организационно-правовая форма и система налогообложения}

Для~реализации проекта выбрана организационно-правовая форма ООО (Общество
с~ограниченной ответственностью) с~системой налогообложения УСН~«Доходы»
(ставка 6~\%). Выбор обусловлен следующими причинами:

\begin{itemize}[nosep]
  \item ООО является предпочтительной формой для~B2B-взаимодействия со~средним и~крупным бизнесом, особенно при~заключении договоров с~государственными или~окологосударственными заказчиками;
  \item ограниченная ответственность участников в~рамках уставного капитала позволяет защитить личное имущество основателей в~случае коммерческих рисков;
  \item УСН~«Доходы» 6~\% оптимальна для~подписочной B2B-модели, поскольку доля расходов в~общей выручке прогнозируется ниже порогового значения 60~\%, при~котором становится выгоднее УСН~«Доходы минус расходы»;
  \item при~достижении объёма выручки и~состава команды, удовлетворяющих требованиям, организация сможет получить IT-аккредитацию и~претендовать на~пониженные тарифы страховых взносов и~налог на~прибыль.
\end{itemize}

Основные коды ОКВЭД, под~которыми будет вестись деятельность, представлены
в~таблице~3.1.3.8.

\tcap{Таблица 3.1.3.8. Коды ОКВЭД.}
{\small
\begin{longtable}{C{2.5cm} L{12.0cm}}
\toprule
Код ОКВЭД & Расшифровка \\
\midrule
\endhead
62.01 & Разработка компьютерного программного обеспечения \\
62.02 & Деятельность консультативная и~работы в~области компьютерных технологий \\
63.11 & Деятельность по~обработке данных, предоставление услуг по~размещению информации и~связанная с~этим деятельность \\
63.11.1 & Деятельность по~созданию и~использованию баз данных и~информационных ресурсов \\
62.09 & Деятельность, связанная с~использованием вычислительной техники и~информационных технологий, прочая \\
72.19 & Научные исследования и~разработки в~области естественных и~технических наук прочие \\
\bottomrule
\end{longtable}
}

\paragraph{Оценка ёмкости рынка (TAM\,--\,SAM\,--\,SOM)}

Целевым рынком является сегмент облачных и~API-сервисов прогнозирования
временных рядов для~российских компаний малого и~среднего бизнеса.
В~таблице~3.1.3.9 представлена многоуровневая оценка ёмкости рынка.

\tcap{Таблица 3.1.3.9. Оценка ёмкости рынка.}
{\small
\begin{longtable}{C{1.2cm} L{10.5cm} L{2.8cm}}
\toprule
Показатель & Описание & Объём \\
\midrule
\endhead
PAM & Potential Available Market~--- мировой рынок облачных ML-сервисов прогнозирования временных рядов (Time Series Forecasting as a Service) & $\approx$~\$5,4 млрд (2025) \\
TAM & Total Addressable Market~--- российский рынок облачных ML-сервисов и~сервисов прогнозирования (по~оценке Strategy Partners и~iKS-Consulting) & $\approx$~12 млрд руб./год \\
SAM & Serviceable Available Market~--- компании МСБ в~РФ в~сегментах ритейла, логистики, энергетики и~телекома, готовые использовать API-сервисы прогнозирования ($\approx$~8000 компаний $\times$ средний чек $\approx$~250~тыс. руб./год) & $\approx$~2,0 млрд руб./год \\
SOM & Serviceable Obtainable Market~--- реально достижимая доля рынка в~горизонте 3~лет с~учётом ресурсов команды (60 активных подписок при~среднем чеке 1,02 млн руб./год) & $\approx$~60 млн руб./год \\
\bottomrule
\end{longtable}
}

По~данным аналитических агентств, общий рынок облачных ML-сервисов в~России
растёт темпом 20\,--\,25~\% в~год благодаря тренду на~цифровизацию
и~импортозамещению зарубежных ML-платформ. Сегмент специализированных сервисов
прогнозирования временных рядов является одной из~наиболее быстрорастущих ниш
внутри этого рынка, что~подтверждает наличие устойчивого спроса и~перспективы
развития проекта.

\paragraph{Экономический эффект для~клиента}

До~внедрения сервиса компания может использовать простые модели
прогнозирования либо содержать собственного специалиста по~анализу данных.
В~обоих случаях сохраняются риски ошибок прогноза, которые приводят
к~избыточным запасам, дефициту товаров, неверному планированию ресурсов
и~дополнительным операционным затратам. Данный подраздел оценивает
экономический эффект для~клиента; потенциальные потери клиента не~включаются
в~расходную часть бюджета проекта.
Актуальность проблемы подтверждается данными IHL~Group: мировой ущерб ритейла
от~дефицита и~избытка товаров оценивается в~1,77~трлн долл. в~2025~году,
а~искажения запасов составляют около 6,5~\% от~глобальных розничных продаж.
Для~расчёта в~работе используется более консервативная оценка потерь~---
3~\% от~оборота. В~качестве расчётного примера рассмотрим компанию
с~ежемесячной выручкой 100~млн руб. Тогда потери от~неточного прогнозирования
составят:
\[ 100\,000\,000 \times 0{,}03 = 3\,000\,000~\text{руб./мес.} \]
Дополнительно для~разработки и~сопровождения собственных моделей клиенту
требуется Data Scientist. По~данным Хабр Карьеры, средняя зарплата такого
специалиста составляет около 251\,666~руб., а~на~hh.ru встречаются вакансии
с~уровнем оплаты 245\,000\,--\,295\,000~руб. Поэтому в~расчёте принимается
значение 250\,000~руб./мес. С~учётом страховых взносов по~ставке 30,2~\%
альтернативные затраты бизнеса составят:
\[ 250\,000 \times 1{,}302 = 325\,500~\text{руб./мес.} \]
Итого совокупные потенциальные потери и~альтернативные расходы клиента
до~внедрения сервиса:
\[ 3\,000\,000 + 325\,500 = 3\,325\,500~\text{руб./мес.} \]
После~внедрения «ГМ-ПВР» компания получает доступ к~гибридному API-сервису
прогнозирования, учитывающему тренд, сезонность, шум и~сложную структуру
временных рядов. Предположим, что~более устойчивый прогноз снижает потери
от~неточного планирования на~25~\%:
\[ 3\,000\,000 \times 0{,}25 = 750\,000~\text{руб./мес.} \]
Кроме того, сервис снижает объём ручных задач по~подбору, обучению
и~сопровождению моделей. Это позволяет не~нанимать дополнительного младшего
специалиста или~перераспределить задачи действующего аналитика. В~расчёте
принимается экономия 150\,000~руб./мес. С~учётом стоимости подписки
85\,000~руб./мес. итоговая экономия составит:
\[ 750\,000 + 150\,000 - 85\,000 = 815\,000~\text{руб./мес. (округлённо).} \]
Годовой экономический эффект на~одного клиента:
\[ 815\,000 \times 12 = 9\,780\,000~\text{руб./год (округлённо).} \]

\paragraph{План продвижения}

Продукт ориентирован на~B2B-рынок компаний малого и~среднего бизнеса.
На~таком рынке ключевыми каналами являются экспертный контент-маркетинг,
прямые продажи и~участие в~отраслевых ML-конференциях. Стандартные методы
B2C-рекламы (таргетированная реклама в~социальных сетях, баннерная реклама)
имеют низкую эффективность, поскольку лицом, принимающим решения, обычно
является руководитель аналитического отдела или~ИТ-директор, которых нужно
выводить на~коммуникацию через~экспертный контент и~кейсы.
Для~повышения эффективности программы продвижения выбран экспертный подход:
систематическая публикация технического и~аналитического контента,
демонстрирующего глубокую экспертизу команды в~области прогнозирования
временных рядов и~гибридных моделей. Примеры таких активностей: ведение
технического блога компании, публикации на~Habr и~vc.ru с~разбором задач
прогнозирования на~сложных профилях рядов, открытые статьи
в~научно-практических изданиях и~доклады на~ML-конференциях. Каждая публикация
содержит нативную привязку к~сервису через~ссылки на~сайт и~демонстрационное
API. Параллельно с~контент-маркетингом ведётся работа по~прямым продажам:
формирование базы потенциальных клиентов по~данным ЕГРЮЛ среди компаний
из~целевых ОКВЭД с~подходящими параметрами по~выручке и~численности персонала;
сегментация базы по~приоритетам; разработка персонализированных предложений
и~серий касаний (e-mail~+ телефонный звонок~+ демонстрация).
В~таблице~3.1.3.10 представлены затраты на~продвижение в~первые три года
реализации проекта. Эти затраты далее включаются в~прогноз прибылей и~убытков
и~движение денежных средств, поэтому план продвижения приводится до~плана
продаж и~финансовых расчётов.

\tcap{Таблица 3.1.3.10. Затраты на~продвижение, тыс. руб.}
{\small
\begin{longtable}{L{9.0cm} C{1.8cm} C{1.8cm} C{1.8cm}}
\toprule
Наименование затрат & Год~1 & Год~2 & Год~3 \\
\midrule
\endhead
Заработная плата контент-менеджера (с~отчислениями) & 180 & 360 & 540 \\
Заработная плата менеджера по~продажам (с~отчислениями) & 0 & 420 & 630 \\
Публикации в~отраслевых изданиях и~на~Habr & 60 & 80 & 100 \\
Разработка и~поддержка корпоративного сайта & 80 & 20 & 20 \\
Участие в~ML-конференциях (Data Fusion, AI~Journey, OpenTalks.AI) & 80 & 120 & 200 \\
Командировки и~пилотные демонстрации у~клиентов & 30 & 80 & 250 \\
Прочие маркетинговые расходы (баннерная реклама на~Habr и~vc.ru) & 20 & 20 & 60 \\
\midrule
\textbf{Итого} & \textbf{450} & \textbf{1100} & \textbf{1800} \\
\bottomrule
\end{longtable}
}

\subsection{Оценка эффективности приложенных решений с~точки зрения затрат и~результатов}

\subsubsection{План продаж}

Монетизация осуществляется через~продажу подписки на~REST~API сервиса.
В~таблице~3.2.1.1 представлен план продаж на~3~года. Расчёт основан
на~консервативной модели роста клиентской базы: в~первый год сервис набирает
клиентов поэтапно, во~второй и~третий годы рост ускоряется за~счёт
сформированной репутации, кейсов внедрения и~партнёрских каналов.

\tcap{Таблица 3.2.1.1. План продаж на~3~года.}
{\small
\begin{longtable}{L{9.0cm} C{1.8cm} C{1.8cm} C{1.8cm}}
\toprule
Показатель & Год~1 & Год~2 & Год~3 \\
\midrule
\endhead
Количество новых подписок за~год, шт. & 12 & 24 & 30 \\
Количество клиентов накопительно, шт. & 12 & 36 & 66 \\
Средняя продолжительность подписки в~год~1, мес. & 6,5 & 12 & 12 \\
Доход от~базовой подписки (85 тыс. руб./мес.), тыс. руб. & 6\,630 & 30\,600 & 56\,100 \\
Доход от~превышения лимитов API, тыс. руб. & 420 & 2\,400 & 5\,200 \\
Доход от~пилотных проектов, тыс. руб. & 1\,500 & 2\,000 & 2\,500 \\
\midrule
\textbf{Итого выручка, тыс. руб.} & \textbf{8\,550} & \textbf{35\,000} & \textbf{63\,800} \\
Операционные расходы, тыс. руб. & 2\,605,1 & 6\,310,2 & 10\,515,3 \\
\textbf{Валовая прибыль, тыс. руб.} & \textbf{5\,944,9} & \textbf{28\,689,8} & \textbf{53\,284,7} \\
\bottomrule
\end{longtable}
}

Операционные расходы включают арендную плату за~облачные GPU-серверы
(масштабируется с~ростом числа клиентов), ФОТ команды поддержки и~развития,
маркетинговые расходы и~налоги. ФОТ команды рассчитан на~основе ранее
определённой стоимости труда участников проекта с~применением ставки страховых
взносов 30,2~\%: в~первый год учитывается одна базовая команда, во~второй~---
две, в~третий~--- три. С~увеличением клиентской базы переменные расходы растут
существенно медленнее выручки за~счёт эффекта масштаба подписочной SaaS-модели.

\subsubsection{Прогноз прибылей и убытков}

На~основании плана продаж и~структуры операционных расходов составлен прогноз
прибылей и~убытков (P\&L) на~трёхлетний горизонт. Результаты представлены
в~таблице~3.2.2.1.

\tcap{Таблица 3.2.2.1. Прогноз прибылей и~убытков, тыс. руб.}
{\small
\begin{longtable}{L{9.0cm} C{1.8cm} C{1.8cm} C{1.8cm}}
\toprule
Показатель & Год~1 & Год~2 & Год~3 \\
\midrule
\endhead
\textbf{Выручка} & \textbf{8\,550} & \textbf{35\,000} & \textbf{63\,800} \\
\quad в~т.\,ч. подписки & 6\,630 & 30\,600 & 56\,100 \\
\quad в~т.\,ч. превышение лимитов & 420 & 2\,400 & 5\,200 \\
\quad в~т.\,ч. пилотные проекты & 1\,500 & 2\,000 & 2\,500 \\
\textbf{Операционные расходы} & \textbf{2\,605,1} & \textbf{6\,310,2} & \textbf{10\,515,3} \\
\quad в~т.\,ч. ФОТ команды (с~отчислениями 30,2~\%) & 1\,505,1 & 3\,010,2 & 4\,515,3 \\
\quad в~т.\,ч. аренда GPU и~облачной инфраструктуры & 450 & 1\,800 & 3\,600 \\
\quad в~т.\,ч. маркетинг и~продвижение & 450 & 1\,100 & 1\,800 \\
\quad в~т.\,ч. прочие операционные расходы & 200 & 400 & 600 \\
Непредвиденные расходы (1~\% от~выручки) & 86 & 350 & 638 \\
\textbf{Прибыль до~налогов} & \textbf{5\,858,9} & \textbf{28\,339,8} & \textbf{52\,646,7} \\
Налог УСН~«Доходы» (6~\% от~выручки) & 513 & 2\,100 & 3\,828 \\
\textbf{Чистая прибыль} & \textbf{5\,345,9} & \textbf{26\,239,8} & \textbf{48\,818,7} \\
\bottomrule
\end{longtable}
}

\tcap{Таблица 3.2.2.2. Таблица денежных потоков, тыс. руб.}
{\small
\begin{longtable}{L{9.0cm} C{1.8cm} C{1.8cm} C{1.8cm}}
\toprule
Статья & Год~1 & Год~2 & Год~3 \\
\midrule
\endhead
Притоки (выручка) & 8\,550 & 35\,000 & 63\,800 \\
Оттоки (операционные расходы~+ непредвиденные~+ налог) & 3\,204,1 & 8\,760,2 & 14\,981,3 \\
Сальдо ДДС за~период & 5\,345,9 & 26\,239,8 & 48\,818,7 \\
Сальдо ДДС накопительно & 5\,345,9 & 31\,585,7 & 80\,404,4 \\
\bottomrule
\end{longtable}
}

\subsubsection{Расчёт показателей окупаемости}

Для~оценки экономической эффективности проекта используются следующие
показатели: чистая приведённая стоимость (NPV), индекс рентабельности (PI),
простой срок окупаемости (PB) и~дисконтированный срок окупаемости (DPP).
В~качестве первоначальных инвестиций принята полная себестоимость разработки
сервиса, рассчитанная в~подразделе~3.1.2.10: $IC = 890\,142{,}8$~руб. Ставка
дисконтирования принята равной 22~\% (как~требуемая доходность инвестора
с~учётом рискового профиля IT-стартапа в~текущих макроэкономических условиях).
Расчёт NPV выполняется по~формуле:
\[ NPV = \sum_{i=1}^{n} \frac{CF_i}{(1+r)^{i}} - IC, \]
где $CF_i$~--- денежный поток $i$-го периода, руб.; $r$~--- ставка
дисконтирования; $IC$~--- первоначальные инвестиции, руб.

\tcap{Таблица 3.2.3.1. Расчёт NPV.}
{\small
\begin{longtable}{C{2.6cm} C{3.2cm} C{4.0cm} C{3.6cm}}
\toprule
Период, год & CF, руб. & Коэф. дисконтирования ($r = 22~\%$) & Текущая стоимость, руб. \\
\midrule
\endhead
0 (инвестиции) & $-890\,142,8$ & 1,0000 & $-890\,142,8$ \\
1 & 5\,345\,900 & 0,8197 & 4\,381\,885,2 \\
2 & 26\,239\,800 & 0,6719 & 17\,629\,535,1 \\
3 & 48\,818\,700 & 0,5507 & 26\,884\,794,3 \\
\midrule
\multicolumn{3}{r}{\textbf{NPV}} & \textbf{48\,006\,071,8 руб.} \\
\bottomrule
\end{longtable}
}

Индекс рентабельности (PI) рассчитывается по~формуле:
\[ PI = \frac{1}{IC}\sum_{i=1}^{n} \frac{CF_i}{(1+r)^{i}} = \frac{NPV + IC}{IC}
   = \frac{48\,006\,071{,}8 + 890\,142{,}8}{890\,142{,}8} \approx 54{,}93. \]
Значение PI существенно превышает~1, что~свидетельствует о~высокой
экономической целесообразности проекта.
Простой срок окупаемости (PB) рассчитывается по~формуле:
\[ PB = \frac{IC}{\overline{CF}}, \]
где $\overline{CF}$~--- средний денежный поток за~период. С~учётом того,
что~денежный поток первого года (5\,345,9~тыс. руб.) уже превышает
первоначальные инвестиции (890,1~тыс. руб.), простой срок окупаемости
составляет:
\[ PB = \frac{890\,142{,}8}{26\,801\,466{,}7} \approx 0{,}03~\text{года.} \]
Дисконтированный срок окупаемости (DPP) рассчитывается по~формуле:
\[ DPP = \min\left\{\, n : \sum_{i=1}^{n} \frac{CF_i}{(1+r)^{i}} \ge IC \,\right\}. \]
Поскольку дисконтированный денежный поток первого года (4\,381\,885,2~руб.)
уже превышает первоначальные инвестиции (890\,142,8~руб.), дисконтированный
срок окупаемости также наступает в~течение первого года:
\[ DPP \approx \frac{890\,142{,}8}{4\,381\,885{,}2} \approx 0{,}20~\text{года.} \]
Стоит отметить, что~в~подразделе~3.2.4 рассматривается более консервативный
сценарий с~подключением одного клиента в~первый месяц и~приростом по~одному
клиенту каждый последующий месяц. По~этому сценарию полная окупаемость затрат
на~разработку достигается к~концу 5-го месяца после~начала коммерческой
деятельности, что~согласуется с~настоящей оценкой DPP в~рамках общего
трёхлетнего бизнес-плана.

\tcap{Таблица 3.2.3.2. Сводные показатели эффективности проекта.}
{\small
\begin{longtable}{L{4.8cm} C{3.4cm} L{6.0cm}}
\toprule
Показатель & Значение & Интерпретация \\
\midrule
\endhead
IC (первоначальные инвестиции) & 890\,142,8 руб. & Полная себестоимость разработки сервиса \\
NPV (чистая приведённая стоимость) & 48\,006\,071,8 руб. & $NPV > 0$~--- проект экономически целесообразен \\
PI (индекс рентабельности) & $\approx 54{,}93$ & Высокая рентабельность вложений \\
PB (простой срок окупаемости) & $\approx 0{,}03$ года & Проект окупается в~течение первого месяца эксплуатации \\
DPP (дисконтированный срок окупаемости) & $\approx 0{,}20$ года & Окупаемость с~учётом ставки 22~\% также в~течение первого квартала \\
Ставка дисконтирования & 22~\% & Принята как~требуемая доходность инвестора \\
\bottomrule
\end{longtable}
}

Расчёты показывают высокую экономическую целесообразность проекта:
положительное значение NPV, существенно превышающий единицу индекс
рентабельности и~короткий срок окупаемости подтверждают возможность возврата
вложений уже на~раннем этапе эксплуатации. Высокая скорость окупаемости
обусловлена особенностями подписочной SaaS-модели: предельные издержки
на~обслуживание дополнительного клиента низкие, а~основные инвестиции уже
понесены на~этапе разработки.

\subsubsection{Срок окупаемости затрат на~разработку (консервативный сценарий)}

Помимо расчёта NPV/PI/DPP, целесообразно рассмотреть консервативный сценарий
выхода на~рынок, в~котором предполагается медленный поэтапный набор клиентской
базы. Этот сценарий показывает нижнюю границу окупаемости и~подтверждает
устойчивость проекта даже в~неблагоприятных условиях. Себестоимость разработки
сервиса, рассчитанная в~подразделе~3.1.2, составляет 890\,142,8~руб. Стоимость
одной месячной подписки на~использование сервиса принимается равной
85\,000~руб. Для~оценки срока окупаемости используется консервативный сценарий
выхода на~рынок: в~первый месяц сервис подключает одного клиента, а~в~каждый
следующий месяц количество новых клиентов увеличивается на~одного. Суммарная
выручка за~$N$~месяцев рассчитывается как~сумма арифметической прогрессии:
\[ S = 85\,000 \cdot \frac{N(N+1)}{2}, \]
где $N$~--- число месяцев с~момента начала продаж. Срок окупаемости
рассчитывается из~условия покрытия затрат на~разработку:
\[ S = 85\,000 \cdot \frac{N(N+1)}{2} \ge 890\,142{,}8. \]
Упростим неравенство:
\[ N(N+1) \ge 20{,}94. \]
Решим данное неравенство подбором целых значений $N$:

\begin{itemize}[nosep]
  \item при $N = 4$: $4 \times 5 = 20 < 20{,}94$, выручка составит $85\,000 \times 20 / 2 = 850\,000$~руб., следовательно, затраты ещё не~окупились;
  \item при $N = 5$: $5 \times 6 = 30 > 20{,}94$, выручка составит $85\,000 \times 30 / 2 = 1\,275\,000$~руб., следовательно, затраты на~разработку окупаются.
\end{itemize}

Таким образом, даже в~консервативном сценарии полная окупаемость затрат
на~исследование, разработку гибридной модели прогнозирования и~создание
REST~API достигается к~концу 5-го месяца после~начала коммерческой
деятельности. Учитывая возможность масштабирования сервиса за~счёт
тиражирования API-доступа без~существенного увеличения затрат на~разработку,
проект обладает высокой инвестиционной привлекательностью.

\subsubsection{Риски проекта}

Для~обеспечения устойчивости проекта проведён анализ потенциальных рисков
с~оценкой вероятности (от~1 до~5) и~силы влияния (от~1 до~5) на~конечный
результат, а~также разработан план мероприятий по~их~нивелированию.
Систематизированные риски и~меры представлены в~таблице~3.2.5.1.

\tcap{Таблица 3.2.5.1. Риски проекта и~меры нивелирования.}
{\footnotesize
\begin{longtable}{L{2.0cm} L{4.2cm} C{0.9cm} C{0.9cm} L{6.0cm}}
\toprule
Категория & Риск & Вер. & Вл. & Меры по~нивелированию \\
\midrule
\endhead
Технологические & Деградация качества прогноза на~специфичных профилях рядов (короткие ряды, часовые данные) & 4 & 4 & Расширение набора декомпозиций (EWT, VMD); автоматический выбор декомпозиции по~характеристикам ряда; явная маркировка неблагоприятных профилей в~API \\
\midrule
Технологические & Сбой в~работе REST~API сервиса & 3 & 4 & Резервирование сервера; автоматический мониторинг; retry-логика на~стороне клиента; SLA с~компенсациями \\
\midrule
Технологические & Рост стоимости GPU-аренды у~облачных провайдеров & 3 & 3 & Параллельная работа с~несколькими провайдерами (Yandex Cloud, VK Cloud, Selectel); оптимизация инференса; кэширование результатов \\
\midrule
Коммерческие & Низкая конверсия из~лидов в~подписку & 4 & 4 & Бесплатный триал-период; интерактивная демонстрационная площадка; кейсы и~отзывы клиентов \\
\midrule
Коммерческие & Появление сильных конкурентов (отечественных или~международных) & 3 & 4 & Постоянное развитие модельной части; патентная защита уникальных решений; «вендор-локин» через~интеграции \\
\midrule
Коммерческие & Самостоятельная разработка решения клиентом (in-house) & 3 & 3 & Демонстрация TCO-преимущества подписки; экспертные кейсы; ускоренный time-to-value за~счёт готового API \\
\midrule
Финансовые & Кассовый разрыв при~медленном наборе клиентской базы & 3 & 4 & Резервный фонд на~6 месяцев; грантовое финансирование Фонда содействия инновациям; выручка от~пилотных проектов \\
\midrule
Финансовые & Задержка оплат от~клиентов & 3 & 3 & Предоплата за~подписку; чёткие договорные условия; автоматическое отключение сервиса при~просрочке свыше 14 дней \\
\midrule
Регуляторные & Ограничения на~трансграничную передачу данных & 2 & 3 & Развёртывание API на~серверах в~РФ; документация по~обработке ПДн; соответствие 152-ФЗ \\
\midrule
Регуляторные & Запрет на~использование open-source библиотек из~недружественных юрисдикций & 2 & 4 & Поддержание форков критичных библиотек; план миграции на~отечественные аналоги \\
\midrule
Кадровые & Уход ключевого ML-инженера & 2 & 5 & Документирование моделей и~кода; передача знаний внутри команды; конкурентная оплата; опционная программа \\
\midrule
Кадровые & Дефицит ML-специалистов на~рынке труда & 3 & 3 & Сотрудничество с~университетами (стажировки, дипломные работы); удалённый формат; внутреннее обучение \\
\bottomrule
\end{longtable}
}

Реализация перечисленных мер по~нивелированию рисков позволит значительно
снизить вероятность их~материализации и~повысить устойчивость проекта на~всех
стадиях коммерциализации. Наиболее существенные риски~--- деградация качества
прогноза на~специфичных профилях рядов и~низкая конверсия из~лидов~---
устраняются за~счёт технического развития сервиса и~продуманной маркетинговой
воронки.

\subsubsection{Вывод по~экономическому разделу}

В~экономическом разделе проведён комплексный анализ коммерческого потенциала
технологического стартапа «ГМ-ПВР»~--- B2B-сервиса краткосрочного
прогнозирования временных рядов через~REST~API на~основе гибридной модели
HybridPlus. Структура раздела приведена в~соответствие с~заданием на~ВКР:
выделены два основных параграфа~--- экономическое обоснование проекта
и~оценка эффективности приложенных решений с~точки зрения затрат и~результатов.
В~рамках подраздела~3.1 выполнены организация и~планирование работ, рассчитана
трудоёмкость девяти этапов разработки (90 чел.-дней разработчика и~25 чел.-дней
руководителя), а~также произведён расчёт полной себестоимости проекта методом
калькулирования. Полная себестоимость составила 890\,142,8~руб., при~этом
наибольшие доли занимают накладные расходы (54,1~\%) и~основная заработная
плата (27,1~\%), что~характерно для~научно-исследовательских проектов в~области
машинного обучения.
В~рамках подраздела~3.1 также разработан комплексный бизнес-план:
сформулированы ценностное предложение и~портрет потребителя B2B-сегмента,
проведены SWOT-анализ и~анализ конкурентов (Amazon Forecast, Azure ML, Yandex
DataSphere, open-source библиотеки), построена бизнес-модель Остервальдера,
определены организационно-правовая форма (ООО) и~система налогообложения
(УСН~«Доходы» 6~\%) с~подбором ОКВЭД. Оценка ёмкости рынка показала следующие
значения: $TAM \approx 12$~млрд руб./год, $SAM \approx 2$~млрд руб./год,
$SOM \approx 60$~млн руб./год в~горизонте 3~лет.
В~рамках подраздела~3.2 сформирован план продаж на~3~года с~выходом на~выручку
63,8~млн руб. в~третий год, рассчитаны прогноз прибылей и~убытков и~таблица
денежных потоков. Расчёт показателей эффективности подтвердил высокую
экономическую целесообразность проекта при~ставке дисконтирования 22~\%:

\begin{itemize}[nosep]
  \item $NPV \approx 48{,}01$~млн руб. ($NPV > 0$~--- проект экономически целесообразен);
  \item $PI \approx 54{,}93$ (значительно больше~1~--- высокая рентабельность);
  \item простой срок окупаемости $PB \approx 0{,}03$~года, то~есть менее одного месяца по~среднему годовому денежному потоку;
  \item дисконтированный срок окупаемости $DPP \approx 0{,}20$~года $\approx 2{,}4$~месяца.
\end{itemize}

Параллельный расчёт в~рамках более консервативного сценария (поэтапное
подключение клиентов с~приростом по~одному в~месяц) показал окупаемость затрат
на~разработку к~концу 5-го месяца коммерческой деятельности, что~согласуется
с~общей оценкой и~подтверждает устойчивость проекта в~неблагоприятных условиях.
Высокая скорость окупаемости обусловлена особенностями подписочной SaaS-модели:
предельные издержки на~обслуживание дополнительного клиента невелики,
а~основные инвестиции концентрируются на~этапе разработки. План продвижения
с~акцентом на~экспертный контент-маркетинг и~прямые продажи приведён
до~финансовых расчётов и~учтён в~P\&L; общий бюджет продвижения на~3~года
составляет 3,35~млн руб. Проведён анализ рисков по~технологической,
коммерческой, финансовой, регуляторной и~кадровой категориям с~разработкой
мер нивелирования.
На~основании проведённого анализа проект «ГМ-ПВР» рекомендуется к~реализации:
он~обладает значимым ценностным предложением для~российских компаний малого
и~среднего бизнеса, экспериментально подтверждённой эффективностью гибридной
MODWT-модели на~сложных профилях рядов, понятной подписочной моделью
монетизации и~высокими расчётными показателями экономической эффективности.
